\documentclass[12pt]{letter}
\usepackage[left=1in,right=1in,top=1in,bottom=1in]{geometry}
\usepackage{lmodern}
\usepackage[T1]{fontenc}
\usepackage{setspace}
\onehalfspacing

\signature{Darcio Genicolo-Martins\\Paulo Furquim de Azevedo\\INSPER, S\~ao Paulo, Brazil}
\address{INSPER\\Rua Quat\'a, 300\\S\~ao Paulo, SP 04546-042\\Brazil}
\date{April 2026}

\begin{document}

\begin{letter}{The Editors\\Journal of Law and Economics\\University of Chicago}

\opening{Dear Editors,}

We submit ``Screening for Bid Rigging with Frequent Losers in Public Procurement'' for consideration at the \textit{Journal of Law and Economics}.

The paper proposes a participation-based triage tool for procurement authorities investigating bid rigging. The screen---\textit{frequent losers} (FL), firms that never win yet bid abnormally often---requires only win/loss records and is deployable wherever participation data exist, including low-capacity enforcement settings where bid-level forensics are infeasible.

Applied to 4.5 million tender-items from S\~ao Paulo's electronic procurement platform (2009--2019), the screen achieves AUC~$= 0.94$ against competition-authority convictions and flags environments with 3.6--7.7\% higher conditional prices. The price association concentrates in competitive markets and where cover bidding is voluntary rather than constraint-driven---consistent with strategic deployment. Five supporting diagnostics are consistent with coordinated cover bidding; the paper's contribution is the screen itself and its external validation.

We believe the paper is well suited for JLE for three reasons:

\begin{enumerate}
\item \textbf{Enforcement triage.} The screen is designed to prioritize suspicious procurement environments for investigation, not to adjudicate firm-level liability. We propose a three-stage enforcement pathway (screen, triage, investigate) and show that the screen's bite is strongest where oversight capacity is weakest---a 12.5$\times$ gradient across purchasing-unit size quartiles---directly relevant to the allocation of scarce enforcement resources.

\item \textbf{Administrability and low data requirements.} The screen runs on participation and outcome records that every procurement system already collects. The triage signal is robust across threshold choices (significant in all 36 IQR $\times$ win-rate cells) and nearly orthogonal to bid-level screens (correlation 0.06), making it a practical first-stage complement to existing detection tools.

\item \textbf{Institutional design implications.} The data are consistent with the view that Brazil's minimum-bidder rule in sealed-bid procurement (convite) creates an incentive for cover bidding. Brazil's recent procurement reform (Lei 14.133/2021), which restructures bidding modalities and minimum-bidder rules, provides a natural laboratory for assessing whether the screen adapts to institutional change.
\end{enumerate}

The manuscript has not been submitted elsewhere. A companion paper developing the full structural model and policy counterfactuals is available as a working paper. All data and code will be made available upon acceptance in compliance with JLE's replication policy.

\closing{Sincerely,}

\end{letter}
\end{document}
